\section{One-time compilation into component profiles}
\label{sec:profiles}

\begin{definition}[Profile census]
For an $E$-orbit $O$, define
\[
 \profile(O)=(|O\cap J_0|,\ldots,|O\cap J_{m-1}|).
\]
The profile census $h_E$ maps a vector $\profile$ to the number of orbits
with that vector. Only positive entries are stored. Write $s=|\supp h_E|$
and $C=\sum_{\profile}h_E(\profile)$.
\end{definition}

A nonempty orbit has a nonzero nonnegative profile. For each atom,
\begin{equation}
 0\le a_j\le\ell_j,
 \qquad
 \sum_{\profile}h_E(\profile)a_j=\ell_j.
 \label{eq:mass}
\end{equation}
The multiplicity $h_E(\profile)$ is binary. A row can therefore describe
$2^{16000}$ orbits without listing them. Equal profiles do not mean that the
orbits are already identified.

\subsection{Capacity-aware scalar weights}

Choose integer radices $R_j>\ell_j$, set
\[
 u_0=1,\qquad u_j=\prod_{i<j}R_i,
 \qquad P=\prod_{j<m}R_j,
\]
and assign scalar weight $u_j$ to every point of $J_j$.

\begin{theorem}[Carry-free profile compilation]
\label{thm:packing}
The sum of these scalar weights on an orbit is
\[
 z(O)=\sum_{j<m}a_j(O)u_j\in[0,P).
\]
Successive quotient--remainder division by $R_0,R_1,\ldots,R_{m-1}$ recovers
$\profile(O)$ uniquely. Consequently one exact scalar weighted orbit
histogram determines the entire profile census, with the same
multiplicities. This holds for every interval relation, including relations
with reflections and static points.
\end{theorem}

\begin{proof}
For any profile satisfying \eqref{eq:mass}, its $j$th coordinate is less than
$R_j$. Reducing $z$ modulo $R_0$ therefore gives $a_0$; subtracting $a_0$ and
dividing by $R_0$ produces the same representation starting with $a_1$.
Induction decodes every coordinate. The usual mixed-radix identity
$\sum_j(R_j-1)u_j=P-1$ bounds $z$. Conversely encoding the decoded digits
reconstructs $z$, so two distinct profiles cannot collide. Summing scalar
weights over an orbit gives exactly the displayed formula regardless of
which pairing descriptions generated that orbit. Grouping equal scalar sums
therefore groups exactly the same orbits as grouping equal profiles.
\end{proof}

Two implementations are supplied. The \texttt{mixed} mode sets
$R_j=\ell_j+1$. The \texttt{binary} mode sets
\[
 b_j=\lceil\log_2(\ell_j+1)\rceil,
 \qquad R_j=2^{b_j},\qquad W=\sum_{j<m}b_j.
\]
Binary mode uses fixed-width nonoverlapping fields, with $W$ total bits.
A decoded field exceeding $\ell_j$ is invalid even if it fits its field.
Both modes reject negative scalar sums and nonzero leftover high digits.

\begin{proposition}[Packing width and its limited optimality]
\label{prop:width}
Mixed-radix packing uses the integer range
$[0,\prod_j(\ell_j+1))$. Any injective code for the \emph{entire} box
$\prod_j\{0,\ldots,\ell_j\}$ needs at least
$\lceil\log_2\prod_j(\ell_j+1)\rceil$ bits in the worst case. Binary mode
uses less than $m$ extra bits beyond the unrounded logarithmic bound, and
$W\le mB$.
\end{proposition}

\begin{proof}
The box has exactly the stated product cardinality. An injective map to
$b$-bit strings needs $2^b$ at least that cardinality. Mixed-radix encoding
is a bijection between the box and its integer range. Rounding each of its
$m$ logarithmic contributions upward adds less than $m$ in total. Finally
$\ell_j\le N$ implies $b_j\le B$.
\end{proof}

This is not an information lower bound for a particular source census,
whose support can be far smaller than the full box. Nor is it a claim that
packing beats vector arithmetic on every machine. A scalar of width $W$
can make multiplication and division more expensive than a short sparse
vector operation. The engineering hypothesis is reduced tuple and dispatch
overhead in some weighted-replay workloads; native measurements are still
required before selecting a default mode.

\begin{example}[Nonuniform capacities]
Take lengths $(2,5,1)$. Mixed mode has radices $(3,6,2)$ and units $(1,3,18)$.
The profile $(2,4,1)$ has code $32$. Decoding gives remainders $2,4,1$.
The full capacity box has $36$ elements and needs six bits. Binary mode has
widths $(2,3,1)$, also six bits. Using the common universe bound $N=8$ for
every field would instead allocate twelve bits. The example illustrates
capacity savings, not a measured native timing claim.
\end{example}

\subsection{Why the compiled table has polynomial encoded size}

Weighted AHT applied to the $m$ basis-vector weight runs gives a compressed
profile output in polynomial time in $k,m,B$. Applying it instead to the
scalar weights above has one coordinate and weight bit sizes $O(W+\log m)$,
again polynomial in $k,m,B$ \cite{aht}. Decoding and coalescing its polynomial
number of output runs yields $h_E$. Thus the fact that $s$ is polynomially
bounded for interval-encoded input follows from classical weighted AHT; it
is not a new sparsity assumption introduced to make the theorem work.

The article deliberately leaves the underlying polynomial as
$P(k,m,B)$ rather than guessing a small exponent for the maintained code.
In particular, a polynomial theorem does not mean that repeated compiled
queries are costless or that one scalar call has constant-size arithmetic.
The software records native statistics without interpreting an AHT cycle
allowance as a bound on all bit operations.

\subsection{Reusing the table for observations}

Let $w_j\in\Z^d$ be a fixed weight vector on atom $J_j$. Define
\[
 A_w(\profile)=\sum_{j<m}a_jw_j.
\]
Then the component-weight histogram is
\begin{equation}
 H_w(v)=\sum_{\profile:A_w(\profile)=v}h_E(\profile).
 \label{eq:observe}
\end{equation}
This formula permits signed weights, including Euler-characteristic charges
when their geometric ownership was previously validated. It costs $O(smd)$
integer multiplications and additions plus grouping, with no orbit search
and no trace transport. Packing is used only for nonnegative source counts;
signed weights are applied after decoding.

More generally, any explicitly computable predicate $F$ on a profile gives
$\sum_{\profile:F(\profile)}h_E(\profile)$. Port incidence is recovered by
checking whether $\sum_{j\in A}a_j$ is positive. Several observations may be
performed together as a vector map, preserving their joint distribution.
Computing separate marginal histograms and later combining them generally
loses correlations; the joint profile table avoids that loss.

\begin{example}[Why marginal caches do not suffice]
A table with profiles $(1,1)$ and $(2,2)$, each of multiplicity one, has the
same two coordinate marginals as the table with profiles $(1,2)$ and $(2,1)$.
Both have atom lengths $(3,3)$. The observation $a_0-a_1$ distinguishes them:
the first gives two zeroes and the second gives $-1$ and $1$. Full joint
profiles, rather than independent atom marginals, are needed for universal
observation reuse.
\end{example}
